\chapter{مفاهیم پایه‌ای}
\label{chap:concepts}

در این فصل، مفاهیم و تعاریف پایه‌ای که در ادامه این پایان‌نامه به آن‌ها نیاز خواهیم داشت معرفی می‌شوند. نخست در \autoref{sec:fl_basics}، اصول یادگیری فدرال و الگوریتم \gls{FedAvg} تشریح می‌گردد. سپس در \autoref{sec:hetero_types}، انواع مختلف ناهمگونی در سیستم‌های توزیع‌شده بررسی می‌شود. در \autoref{sec:byzantine}، مفهوم خطای بیزانتین و تهدیدات امنیتی در یادگیری فدرال معرفی می‌شود. در \autoref{sec:cloud_concepts}، ویژگی‌های محیط‌های ابری چندگانه که زمینه اصلی این پژوهش هستند بیان می‌شود و در نهایت در \autoref{sec:trust_concepts}، مفاهیم اعتماد و اجماع در سیستم‌های توزیع‌شده مطرح می‌گردد.

\section{یادگیری فدرال}
\label{sec:fl_basics}

\subsection{تعریف و اصول کلی}

\begin{ntdefinition}
یادگیری فدرال یک رویکرد یادگیری ماشین توزیع‌شده است که در آن، یک مدل سراسری به‌صورت مشارکتی میان چندین گره آموزش می‌بیند، بدون آنکه داده‌های محلی هر گره از محیط آن خارج شود.
\end{ntdefinition}

در فرمول‌بندی کلی یادگیری فدرال، $N$ گره داریم که هر کدام مجموعه داده محلی $\mathcal{D}_i$ دارند. هدف کمینه‌سازی تابع زیر است:

\begin{equation}
    F(w) = \sum_{i=1}^{N} \frac{|\mathcal{D}_i|}{|\mathcal{D}|} F_i(w)
    \label{eq:fed_objective}
\end{equation}

که در آن $F_i(w) = \frac{1}{|\mathcal{D}_i|}\sum_{x \in \mathcal{D}_i} \ell(w; x)$ تابع خطای محلی گره $i$ است و $\ell(w; x)$ خطای پیش‌بینی مدل با پارامترهای $w$ روی نمونه $x$ است.

\begin{note}
    تفاوت اساسی یادگیری فدرال با یادگیری توزیع‌شده سنتی در این است که در یادگیری فدرال، توزیع داده میان گره‌ها می‌تواند کاملاً ناهمگون باشد و تعداد نمونه‌های هر گره نیز متفاوت است. در یادگیری توزیع‌شده سنتی، معمولاً داده به‌صورت مساوی بین گره‌ها تقسیم می‌شود.
\end{note}

\subsection{الگوریتم \lr{FedAvg}}

الگوریتم \lr{FedAvg} که توسط \lr{McMahan} و همکاران \cite{mcmahan2017communication} معرفی شد، پایه‌ای‌ترین روش برای حل مسئله \gls{FL} است. این الگوریتم در هر دور ارتباطی $t$ مراحل زیر را طی می‌کند:

\begin{enumerate}
    \item سرور زیرمجموعه‌ای از گره‌ها $S_t \subseteq [N]$ را با اندازه $K$ انتخاب می‌کند.
    \item سرور مدل سراسری $w^t$ را به همه گره‌های انتخاب‌شده ارسال می‌کند.
    \item هر گره $i \in S_t$ برای $E$ دوره محلی، \gls{SGD} را روی داده‌های خود اجرا می‌کند:
    \begin{equation}
        w_i^{t+1} \leftarrow w_i^t - \eta \nabla F_i(w_i^t)
        \label{eq:local_sgd}
    \end{equation}
    \item هر گره به‌روزرسانی $\Delta_i^t = w_i^{t+1} - w^t$ را به سرور ارسال می‌کند.
    \item سرور با میانگین وزن‌دار، مدل را به‌روز می‌کند:
    \begin{equation}
        w^{t+1} = w^t + \sum_{i \in S_t} \frac{|\mathcal{D}_i|}{\sum_{j \in S_t}|\mathcal{D}_j|} \Delta_i^t
        \label{eq:fedavg}
    \end{equation}
\end{enumerate}

\begin{warning}{محدودیت‌های \lr{FedAvg}}
    الگوریتم \lr{FedAvg} تحت فرضیات زیر طراحی شده است که در محیط‌های ابری واقعی اغلب برقرار نیستند: (الف) همه گره‌ها صادق و قابل‌اعتماد هستند، (ب) داده‌ها تقریباً به‌صورت \gls{IID} توزیع شده‌اند، (ج) همه گره‌ها در زمان مشابهی پاسخ می‌دهند.
\end{warning}

\subsection{همگرایی یادگیری فدرال}

تحلیل همگرایی \lr{FedAvg} در حالت عمومی پیچیده است. برای حالتی که تابع $F$ $\mu$-محدب قوی و $L$-لیپشیتز پیوسته است، می‌توان نشان داد که خطای همگرایی از مرتبه زیر است:

\begin{equation}
    \mathbb{E}[F(w^T)] - F(w^*) \leq \mathcal{O}\left(\frac{1}{T} + \Gamma\right)
    \label{eq:convergence}
\end{equation}

که در آن $\Gamma$ یک اصطلاح واگرایی است که به تفاوت توزیع داده‌ها میان گره‌ها بستگی دارد. در توزیع \gls{IID}، $\Gamma = 0$ است، اما در توزیع \gls{NonIID}، این اصطلاح می‌تواند بزرگ باشد و همگرایی را کُند یا حتی متوقف کند.


\section{انواع ناهمگونی در سیستم‌های توزیع‌شده}
\label{sec:hetero_types}

\subsection{ناهمگونی محاسباتی}

ناهمگونی محاسباتی به تفاوت در توان پردازشی گره‌های مختلف اشاره دارد. در یک محیط ابری، این تفاوت می‌تواند از چند برابر تا چند ده برابر باشد.

\begin{ntdefinition}
سرعت محاسباتی نسبی گره $i$ را با $s_i \in (0, 1]$ نشان می‌دهیم که در آن $s_i = 1$ بیانگر سریع‌ترین گره و $s_i \to 0$ بیانگر کُندترین گره ممکن است.
\end{ntdefinition}

اگر زمان انجام یک دوره آموزش محلی توسط سریع‌ترین گره برابر $T_0$ باشد، آنگاه زمان مورد نیاز گره $i$ برای همان تعداد دوره، برابر $T_0 / s_i$ است. در یک محیط ابری واقعی، $s_i$ می‌تواند به سختی‌افزار نمونه ابری (وجود \gls{GPU}، تعداد هسته‌های \gls{CPU}، حافظه)، بار فعلی نمونه (وجود فرایندهای موازی دیگر) و کارایی کد آموزشی محلی بستگی داشته باشد.

\subsection{ناهمگونی شبکه}

ناهمگونی شبکه به تفاوت در کیفیت کانال ارتباطی میان گره‌ها و سرور اشاره دارد. این ناهمگونی به سه پارامتر اصلی زیر وابسته است:

\begin{itemize}
    \tick \textbf{تأخیر (\lr{Latency})}: زمانی که یک بسته داده برای رفتن از گره به سرور (یا بالعکس) طی می‌کند. در محیط‌های ابری، این مقدار از کمتر از یک میلی‌ثانیه (درون یک ناحیه) تا بیش از ۱۵۰ میلی‌ثانیه (بین دو قاره) متفاوت است.
    
    \tick \textbf{پهنای باند}: سرعت انتقال داده. ارسال یک مدل با میلیون‌ها پارامتر در یک پیوند با پهنای باند کم می‌تواند چندین دقیقه طول بکشد.
    
    \tick \textbf{نرخ افت بسته}: احتمال از دست رفتن یک بسته در مسیر. در محیط‌های \gls{WAN}، این احتمال ناچیز نیست.
\end{itemize}

\begin{ntdefinition}
گره کُند (\gls{Straggler}) به گره‌ای گفته می‌شود که به دلیل ترکیبی از محدودیت‌های محاسباتی یا شبکه‌ای، در یک دور ارتباطی معین نمی‌تواند پیش از پایان یافتن زمان مجاز، پاسخ خود را به سرور ارسال کند.
\end{ntdefinition}

\subsection{ناهمگونی داده}

ناهمگونی داده مهم‌ترین و پیچیده‌ترین نوع ناهمگونی در یادگیری فدرال است. این ناهمگونی ناشی از آن است که داده‌های موجود در گره‌های مختلف توزیع آماری یکسانی ندارند.

\begin{ntdefinition}
داده‌ها را \gls{IID} می‌نامیم اگر نمونه‌های موجود در هر گره از یک توزیع مشترک و مستقل از یکدیگر نمونه‌برداری شده باشند. در غیر این صورت، داده‌ها \gls{NonIID} هستند.
\end{ntdefinition}

\subsubsection{مدل‌سازی ناهمگونی با توزیع دیریکله}

یکی از روش‌های رایج برای مدل‌سازی ناهمگونی داده، استفاده از توزیع دیریکله است. در این روش، بردار احتمال توزیع برچسب هر گره از توزیع دیریکله با پارامتر $\alpha$ نمونه‌برداری می‌شود:

\begin{equation}
    p_i \sim \text{Dir}(\alpha \cdot \mathbf{1}_C)
    \label{eq:dirichlet}
\end{equation}

که در آن $C$ تعداد کلاس‌ها و $\mathbf{1}_C$ بردار یک با طول $C$ است. پارامتر $\alpha$ میزان ناهمگونی را کنترل می‌کند:

\begin{itemize}
    \tick هنگامی که $\alpha \to \infty$، توزیع به حالت \gls{IID} نزدیک می‌شود.
    \tick هنگامی که $\alpha = 1$، ناهمگونی متوسطی داریم.
    \tick هنگامی که $\alpha \to 0$، هر گره فقط نمونه‌های یک یا چند کلاس را دارد (ناهمگونی شدید).
\end{itemize}


\section{خطای بیزانتین و تهدیدات امنیتی}
\label{sec:byzantine}

\subsection{تعریف خطای بیزانتین}

\begin{ntdefinition}
در یک سیستم توزیع‌شده، یک گره بیزانتین (\gls{Byzantine}) به گره‌ای گفته می‌شود که رفتار دلخواه و غیرقابل پیش‌بینی دارد. این رفتار ممکن است شامل ارسال مقادیر تصادفی، ارسال مقادیر حساب‌شده برای آسیب‌رسانی، یا عدم پاسخ‌دهی باشد.
\end{ntdefinition}

در زمینه یادگیری فدرال، یک گره بیزانتین می‌تواند به‌روزرسانی‌هایی ارسال کند که مدل سراسری را به سمت نادرستی سوق دهد. مهم‌ترین انواع حملات عبارتند از:

\subsubsection{حمله نویز تصادفی}
در این حمله، گره مخرب به‌روزرسانی واقعی خود را با نویز گاوسی ترکیب می‌کند:
\begin{equation}
    \tilde{\Delta}_i = \Delta_i + \mathcal{N}(0, \sigma^2 I)
    \label{eq:noise_attack}
\end{equation}
این حمله ضعیف‌ترین نوع حمله است، چون معمولاً روش‌های ساده مبتنی بر اندازه نرم می‌توانند آن را شناسایی کنند.

\subsubsection{حمله جهت‌معکوس‌سازی}
در این حمله که به آن \lr{Sign-Flip} نیز گفته می‌شود، گره مخرب علامت به‌روزرسانی خود را معکوس می‌کند:
\begin{equation}
    \tilde{\Delta}_i = -\lambda \Delta_i, \quad \lambda > 0
    \label{eq:sign_flip}
\end{equation}
این حمله خطرناک‌تر است چون به‌روزرسانی مخرب ممکن است اندازه‌ای مشابه به‌روزرسانی‌های صادق داشته باشد.

\subsubsection{حمله مقیاس‌افزایی}
در این حمله، گره مخرب به‌روزرسانی خود را با یک ضریب بزرگ مقیاس می‌دهد:
\begin{equation}
    \tilde{\Delta}_i = \kappa \Delta_i, \quad \kappa \gg 1
    \label{eq:scaling_attack}
\end{equation}
هدف این حمله آن است که تأثیر گره مخرب در تجمیع بسیار بزرگ‌تر از گره‌های صادق باشد.

\subsection{روش‌های موجود برای مقاوم‌سازی}

\subsubsection{الگوریتم \lr{Krum}}
\lr{Krum} \cite{blanchard2017machine} یک الگوریتم تجمیع مقاوم است که در هر دور، $m$ گره را انتخاب می‌کند. برای هر گره $i$، امتیاز \lr{Krum} به صورت زیر تعریف می‌شود:

\begin{equation}
    s_i^{Krum} = \sum_{j \in \mathcal{N}_i} \|\Delta_i - \Delta_j\|^2
    \label{eq:krum_score}
\end{equation}

که در آن $\mathcal{N}_i$ مجموعه $n-f-2$ نزدیک‌ترین همسایه گره $i$ است. گره‌هایی با کمترین امتیاز انتخاب می‌شوند.

\begin{nttheorem}
\label{theorem:krum}
اگر $n \geq 2f + 3$ باشد که در آن $f$ تعداد گره‌های بیزانتین است، آنگاه \lr{Krum} خاصیت $(\alpha, f)$-بیزانتین-مقاومت را دارد.
\end{nttheorem}

\begin{proof}
    اثبات کامل در \cite{blanchard2017machine} آمده است. ایده اصلی اثبات این است که فاصله میان به‌روزرسانی‌های گره‌های صادق محدود است، در حالی که گره‌های بیزانتین می‌توانند به‌روزرسانی‌هایی ارسال کنند که از این مجموعه فاصله زیادی دارند. بنابراین گره‌های صادق امتیاز \lr{Krum} پایین‌تری خواهند داشت.
\end{proof}

\subsubsection{میانگین کوتاه‌شده}
میانگین کوتاه‌شده \cite{yin2018byzantine} یک روش تجمیع مختصه‌ای است که در هر مختصه، $\beta$ درصد از بزرگ‌ترین و کوچک‌ترین مقادیر را حذف می‌کند:

\begin{equation}
    [\text{TrimMean}(\Delta_1, \ldots, \Delta_n)]_k = \frac{1}{n - 2\lfloor\beta n\rfloor} \sum_{i=\lfloor\beta n\rfloor + 1}^{n - \lfloor\beta n\rfloor} \Delta_{\sigma_k(i), k}
    \label{eq:trimmed_mean}
\end{equation}

که در آن $\sigma_k$ مرتب‌سازی مختصه $k$ام به‌روزرسانی‌هاست.


\section{محیط‌های ابری چندگانه}
\label{sec:cloud_concepts}

\subsection{تعریف و ویژگی‌ها}

\begin{ntdefinition}
محیط ابری چندگانه (\lr{Multi-Cloud}) به وضعیتی اطلاق می‌شود که در آن یک سازمان از خدمات بیش از یک ارائه‌دهنده ابری به‌صورت هم‌زمان استفاده می‌کند. این رویکرد با هدف جلوگیری از وابستگی به یک ارائه‌دهنده، بهینه‌سازی هزینه‌ها و بهره‌برداری از بهترین قابلیت‌های هر ارائه‌دهنده اتخاذ می‌شود.
\end{ntdefinition}

در زمینه یادگیری فدرال، محیط ابری چندگانه به معنای آن است که گره‌های شرکت‌کننده در فرایند یادگیری می‌توانند در نواحی مختلف یک ارائه‌دهنده یا روی ارائه‌دهندگان مختلف مستقر باشند.

\subsection{چالش‌های شبکه‌ای در محیط ابری چندگانه}

ارتباط شبکه‌ای در یک محیط ابری چندگانه چند ویژگی اساسی دارد که آن را از شبکه‌های محلی متمایز می‌کند:

\begin{itemize}
    \tick ارتباط بین دو نمونه در یک ناحیه معمولاً تأخیر بسیار کمی دارد (زیر یک میلی‌ثانیه)، اما ارتباط بین دو ناحیه مختلف می‌تواند تأخیر چند ده میلی‌ثانیه داشته باشد.
    
    \tick ارتباط بین دو ارائه‌دهنده ابری مختلف از طریق اینترنت عمومی برقرار می‌شود که تأخیر بیشتری دارد و ضمانت کیفیت سرویس (\lr{QoS}) کمتری ارائه می‌دهد.
    
    \tick هزینه انتقال داده بین ارائه‌دهندگان ابری مختلف (خروجی) می‌تواند قابل‌توجه باشد.
\end{itemize}

\begin{info}
    در این پژوهش، گره‌های شرکت‌کننده در یادگیری فدرال معادل نمونه‌های ابری (ماشین‌های مجازی یا کانتینرها) هستند که در نواحی مختلف یک یا چند ارائه‌دهنده ابری مستقر هستند. ناهمگونی شبکه‌ای مدل‌شده در این پژوهش، بازتاب واقعیت ارتباطات درون‌ناحیه‌ای، بین‌ناحیه‌ای و بین‌ارائه‌دهنده‌ای در چنین محیط‌هایی است.
\end{info}


\section{مفاهیم اعتماد و اجماع}
\label{sec:trust_concepts}

\subsection{اجماع در سیستم‌های توزیع‌شده}

\begin{ntdefinition}
مسئله اجماع در سیستم‌های توزیع‌شده این است که چگونه گره‌های مختلف می‌توانند در مورد یک مقدار یا تصمیم مشترک به توافق برسند، حتی در حضور خرابی‌ها یا رفتارهای مخرب.
\end{ntdefinition}

در یادگیری فدرال، مسئله تجمیع مدل در واقع یک نوع مسئله اجماع است: سرور باید بر اساس ورودی‌های دریافتی از گره‌های مختلف (که برخی ممکن است معیوب یا مخرب باشند) به یک مدل مشترک برسد.

\subsection{مفهوم اعتماد در سیستم‌های توزیع‌شده}

\begin{ntdefinition}
امتیاز اعتماد یک گره، یک کمیت عددی است که رفتار تاریخی آن گره را در یک سیستم توزیع‌شده خلاصه می‌کند. هر چه این امتیاز بالاتر باشد، گره مورد نظر در گذشته رفتار قابل‌اطمینان‌تری داشته و سهم بیشتری در تصمیم‌گیری مشترک دریافت می‌کند.
\end{ntdefinition}

استفاده از امتیاز اعتماد در سیستم‌های توزیع‌شده چند مزیت اساسی دارد. نخست، یک خرابی گذرا (مانند یک بسته از دست رفته یا یک تأخیر مقطعی) موجب محرومیت دائمی گره نمی‌شود. دوم، گره‌هایی که به‌طور مستمر رفتار غیرعادی دارند به‌تدریج تأثیر کمتری در سیستم خواهند داشت. سوم، سیستم می‌تواند به تغییر رفتار گره‌ها در طول زمان واکنش نشان دهد.

\subsection{میانگین متحرک نمایی}

\begin{ntdefinition}
میانگین متحرک نمایی (\gls{EMA}) یک روش برای محاسبه میانگین وزن‌دار دنباله‌ای از مقادیر است که در آن مقادیر اخیر وزن بیشتری دارند:
\begin{equation}
    \bar{x}^t = \mu \bar{x}^{t-1} + (1 - \mu) x^t
    \label{eq:ema}
\end{equation}
که در آن $\mu \in [0, 1)$ پارامتر اینرسی است. هر چه $\mu$ به یک نزدیک‌تر باشد، تغییرات آهسته‌تر است.
\end{ntdefinition}

\gls{EMA} در روش \gls{TWAC} پیشنهادی به دو منظور استفاده می‌شود: نخست برای به‌روزرسانی امتیازات اعتماد گره‌ها به گونه‌ای که تغییرات ناگهانی صاف شوند، و دوم برای نگهداری یک جهت مرجع که نمایانگر میانگین وزن‌دار جهت به‌روزرسانی‌های تجمیع‌شده گذشته است.

\begin{ntremember}
    مهم‌ترین مفاهیم این فصل که در ادامه مکرراً استفاده خواهند شد: (الف) فرمول تجمیع \lr{FedAvg} در معادله \eqref{eq:fedavg}، (ب) مدل ناهمگونی داده با توزیع دیریکله در معادله \eqref{eq:dirichlet}، (ج) تعریف انواع حملات بیزانتین در معادلات \eqref{eq:noise_attack} تا \eqref{eq:scaling_attack}، و (د) تعریف \gls{EMA} در معادله \eqref{eq:ema}.
\end{ntremember}
